\documentclass[11pt,a4paper]{article}
\usepackage[utf8]{inputenc}
\usepackage[spanish]{babel}
\usepackage{amsmath}
\usepackage{amsfonts}
\usepackage{amssymb}
\usepackage{geometry}
\geometry{top=2.5cm, bottom=2.5cm, left=2.5cm, right=2.5cm}

\begin{document}

\begin{center}
    \textbf{\Large SOLUCIONARIO OFICIAL} \\[0.2cm]
    \textbf{\large Exámenes Finales - Versiones A, B, C} \\[0.4cm]
\end{center}

\hrule
\vspace{0.4cm}
\subsection*{Examen Versión A}\textbf{I. Teoría:}\begin{enumerate}    \item \textbf{Respuesta:} Permanecerá en reposo o en movimiento rectilíneo uniforme. \\ \textit{De acuerdo con la Primera Ley de Newton, si la fuerza neta que actúa sobre un cuerpo es nula, el cuerpo necesariamente:}
    \item \textbf{Respuesta:} Vectorial. \\ \textit{La fuerza es una magnitud física de carácter:}
    \item \textbf{Respuesta:} Masa. \\ \textit{La propiedad física de un cuerpo que mide cuantitativamente su inercia (resistencia al cambio en su estado de movimiento) es la:}
    \item \textbf{Respuesta:} Es la fuerza de atracción gravitatoria que ejerce la Tierra sobre él ($w = m \cdot g$). \\ \textit{A diferencia de la masa, el peso de un objeto:}
    \item \textbf{Respuesta:} El coeficiente de fricción cinética ($\mu_c$) es siempre menor que el coeficiente de fricción estática ($\mu_e$). \\ \textit{Con respecto a la fuerza de fricción entre dos superficies secas en contacto, se cumple que:}
    \item \textbf{Respuesta:} La pendiente de la recta tangente a la curva en ese punto. \\ \textit{En un gráfico de posición en función del tiempo ($x$ vs $t$), la velocidad instantánea en un tiempo determinado equivale a:}
    \item \textbf{Respuesta:} $45^\circ$ \\ \textit{Despreciando la resistencia del aire, el ángulo de disparo respecto a la horizontal que produce el máximo alcance horizontal para un proyectil es:}
    \item \textbf{Respuesta:} Directamente proporcional a la fuerza neta e inversamente proporcional a su masa. \\ \textit{La Segunda Ley de Newton establece matemáticamente que la aceleración de un cuerpo es:}
\end{enumerate}\textbf{II. Problemas Prácticos:}\begin{enumerate}    \item \textbf{Respuesta:} $12.5 \text{ m}$ \\ \textit{Un estudiante deja caer un tomate (desde el reposo) desde la azotea de un edificio y mide que tarda exactamente $1.60 \text{ s}$ en golpear el suelo. Despreciando la resistencia del aire y usando $g = 9.80 \text{ m/s}^2$, calcule la altura del edificio:}
    \item \textbf{Respuesta:} $10.0 \text{ N}$ \\ \textit{Determine la magnitud de la fuerza neta sobre una partícula que recibe la acción de las fuerzas $\vec{F}_1 = (5.00 \text{ N}, 4.00 \text{ N})$ y $\vec{F}_2 = (3.00 \text{ N}, 2.00 \text{ N})$:}
    \item \textbf{Respuesta:} $100 \text{ rad/s}$ \\ \textit{Un estudiante mide que las aspas de un abanico de techo realizan exactamente $40 \text{ vueltas}$ en un intervalo de $2.50 \text{ s}$. Calcule la velocidad angular $\omega$ del abanico en rad/s (utilice $\pi \approx 3.14$):}
    \item \textbf{Respuesta:} $28.4 \text{ m}$ \\ \textit{Un automóvil que viaja en línea recta a $50.0 \text{ km/h}$ ($13.9 \text{ m/s}$) pisa los frenos durante $3.00 \text{ s}$ reduciendo su velocidad de forma uniforme hasta $18.0 \text{ km/h}$ ($5.00 \text{ m/s}$). Calcule la distancia recorrida ($\Delta x$) durante el frenado:}
\end{enumerate}\vspace{0.4cm} \hrule \vspace{0.4cm}\subsection*{Examen Versión B}\textbf{I. Teoría:}\begin{enumerate}    \item \textbf{Respuesta:} Permanecerá en reposo o en movimiento rectilíneo uniforme. \\ \textit{De acuerdo con la Primera Ley de Newton, si la fuerza neta que actúa sobre un cuerpo es nula, el cuerpo necesariamente:}
    \item \textbf{Respuesta:} El cuerpo B ejerce sobre A una fuerza de igual magnitud, misma dirección y sentido contrario ($\vec{F}_{BA} = -\vec{F}_{AB}$). \\ \textit{De acuerdo con la Tercera Ley de Newton (Acción-Reacción), si un cuerpo A ejerce una fuerza $\vec{F}_{AB}$ sobre un cuerpo B, entonces:}
    \item \textbf{Respuesta:} $45^\circ$ \\ \textit{Despreciando la resistencia del aire, el ángulo de disparo respecto a la horizontal que produce el máximo alcance horizontal para un proyectil es:}
    \item \textbf{Respuesta:} Masa. \\ \textit{La propiedad física de un cuerpo que mide cuantitativamente su inercia (resistencia al cambio en su estado de movimiento) es la:}
    \item \textbf{Respuesta:} El coeficiente de fricción cinética ($\mu_c$) es siempre menor que el coeficiente de fricción estática ($\mu_e$). \\ \textit{Con respecto a la fuerza de fricción entre dos superficies secas en contacto, se cumple que:}
    \item \textbf{Respuesta:} Vectorial. \\ \textit{La fuerza es una magnitud física de carácter:}
    \item \textbf{Respuesta:} Es la fuerza de atracción gravitatoria que ejerce la Tierra sobre él ($w = m \cdot g$). \\ \textit{A diferencia de la masa, el peso de un objeto:}
    \item \textbf{Respuesta:} Directamente proporcional a la fuerza neta e inversamente proporcional a su masa. \\ \textit{La Segunda Ley de Newton establece matemáticamente que la aceleración de un cuerpo es:}
\end{enumerate}\textbf{II. Problemas Prácticos:}\begin{enumerate}    \item \textbf{Respuesta:} $2.43 \text{ N}$ \\ \textit{Un cuerpo pesa $14.7 \text{ N}$ en la Tierra ($g = 9.80 \text{ m/s}^2$). ¿Cuánto pesará este mismo cuerpo en la superficie de la Luna, donde la aceleración de la gravedad es de $1.62 \text{ m/s}^2$?:}
    \item \textbf{Respuesta:} $-6.60 \text{ N}$ \\ \textit{Un cuerpo de masa $m = 3.00 \text{ kg}$ acelera horizontalmente a razón de $a_x = +0.800 \text{ m/s}^2$ (hacia la derecha). Si actúan dos fuerzas sobre el eje $x$ y una de ellas es $F_{2x} = +9.00 \text{ N}$, determine la componente de la otra fuerza $F_{1x}$:}
    \item \textbf{Respuesta:} $-1.60 \text{ m/s}^2$ \\ \textit{Una bola de golf se desliza por el césped con una velocidad inicial de $6.00 \text{ m/s}$ y llega al hoyo a $2.00 \text{ m/s}$ después de recorrer una distancia de $10.0 \text{ m}$ en línea recta. Determine su aceleración constante:}
    \item \textbf{Respuesta:} $12.5 \text{ m}$ \\ \textit{Un estudiante deja caer un tomate (desde el reposo) desde la azotea de un edificio y mide que tarda exactamente $1.60 \text{ s}$ en golpear el suelo. Despreciando la resistencia del aire y usando $g = 9.80 \text{ m/s}^2$, calcule la altura del edificio:}
\end{enumerate}\vspace{0.4cm} \hrule \vspace{0.4cm}\subsection*{Examen Versión C}\textbf{I. Teoría:}\begin{enumerate}    \item \textbf{Respuesta:} Vectorial. \\ \textit{La fuerza es una magnitud física de carácter:}
    \item \textbf{Respuesta:} Directamente proporcional a la fuerza neta e inversamente proporcional a su masa. \\ \textit{La Segunda Ley de Newton establece matemáticamente que la aceleración de un cuerpo es:}
    \item \textbf{Respuesta:} El coeficiente de fricción cinética ($\mu_c$) es siempre menor que el coeficiente de fricción estática ($\mu_e$). \\ \textit{Con respecto a la fuerza de fricción entre dos superficies secas en contacto, se cumple que:}
    \item \textbf{Respuesta:} $45^\circ$ \\ \textit{Despreciando la resistencia del aire, el ángulo de disparo respecto a la horizontal que produce el máximo alcance horizontal para un proyectil es:}
    \item \textbf{Respuesta:} El centro de la trayectoria circular. \\ \textit{En un Movimiento Circular Uniforme (MCU), aunque la rapidez es constante, la velocidad cambia de dirección, lo que implica una aceleración centrípeta dirigida hacia:}
    \item \textbf{Respuesta:} Permanecerá en reposo o en movimiento rectilíneo uniforme. \\ \textit{De acuerdo con la Primera Ley de Newton, si la fuerza neta que actúa sobre un cuerpo es nula, el cuerpo necesariamente:}
    \item \textbf{Respuesta:} Masa. \\ \textit{La propiedad física de un cuerpo que mide cuantitativamente su inercia (resistencia al cambio en su estado de movimiento) es la:}
    \item \textbf{Respuesta:} El cuerpo B ejerce sobre A una fuerza de igual magnitud, misma dirección y sentido contrario ($\vec{F}_{BA} = -\vec{F}_{AB}$). \\ \textit{De acuerdo con la Tercera Ley de Newton (Acción-Reacción), si un cuerpo A ejerce una fuerza $\vec{F}_{AB}$ sobre un cuerpo B, entonces:}
\end{enumerate}\textbf{II. Problemas Prácticos:}\begin{enumerate}    \item \textbf{Respuesta:} $0.262 \text{ m/s}$ \\ \textit{Una cápsula de pasajeros en una rueda de la fortuna de $150 \text{ m}$ de diámetro gira a rapidez constante y le toma exactamente $30.0 \text{ minutos}$ ($1800 \text{ s}$) completar una vuelta entera. Calcule su rapidez lineal:}
    \item \textbf{Respuesta:} $-6.60 \text{ N}$ \\ \textit{Un cuerpo de masa $m = 3.00 \text{ kg}$ acelera horizontalmente a razón de $a_x = +0.800 \text{ m/s}^2$ (hacia la derecha). Si actúan dos fuerzas sobre el eje $x$ y una de ellas es $F_{2x} = +9.00 \text{ N}$, determine la componente de la otra fuerza $F_{1x}$:}
    \item \textbf{Respuesta:} $100 \text{ rad/s}$ \\ \textit{Un estudiante mide que las aspas de un abanico de techo realizan exactamente $40 \text{ vueltas}$ en un intervalo de $2.50 \text{ s}$. Calcule la velocidad angular $\omega$ del abanico en rad/s (utilice $\pi \approx 3.14$):}
    \item \textbf{Respuesta:} $10.0 \text{ N}$ \\ \textit{Determine la magnitud de la fuerza neta sobre una partícula que recibe la acción de las fuerzas $\vec{F}_1 = (5.00 \text{ N}, 4.00 \text{ N})$ y $\vec{F}_2 = (3.00 \text{ N}, 2.00 \text{ N})$:}
\end{enumerate}\vspace{0.4cm} \hrule \vspace{0.4cm}\end{document}